\section{Finite search and mathematical structure}\label{sec:research}
Qiushi Engine's investigation connected changes in geometric
representation with finite optimization and counting problems.
For replacements, the objective is the net gain after all required
deletions. For lattice sections, the relevant object is a linear
statistic of fibre multiplicities. These reductions retain the
constraints that determine the final point count.

\subsection{An exact exchange objective}
Let $X\subset S^{d-1}$ be a kissing configuration and let
$C\subset S^{d-1}\setminus X$ be a finite candidate set.
For $c\in C$, write
$F(c)=\{x\in X:\ip{c}{x}>1/2\}$. Join two candidates when they
are incompatible. For an independent set $I\subset C$,
\begin{equation}\label{eq:general-exchange}
 X_I=(X\setminus F(I))\cup I,\qquad
 F(I)=\bigcup_{c\in I}F(c),
 \qquad |X_I|-|X|=|I|-|F(I)|.
\end{equation}
The retained old points are compatible, every old--new conflict
has been removed, and the selected new points are pairwise compatible.
Thus maximizing this expression over a candidate pool is an exact
finite subproblem of the geometric search.

For binary variables $z_c$ and $r_x$ indicating insertion and
removal, respectively, the corresponding integer program is
\[
 \max\left(\sum_{c\in C}z_c-\sum_{x\in X}r_x\right),\qquad
 z_c+z_{c'}\le1\ (cc'\in E),\quad
 r_x\ge z_c\ (x\in F(c)).
\]
At an optimum, $r_x=1$ precisely on the union of the selected
conflict sets. When candidates and deletions are disjoint bundles
of internally compatible points, their cardinalities become the
coefficients in the objective.
This is the weighted form behind a support exchange, a signed
replacement, and a first/second-layer head selection.

Write $g(I)=|I|-|F(I)|$. If $c\notin I$ is compatible with $I$, then
\[
 g(I\cup\{c\})-g(I)=1-|F(c)\setminus F(I)|.
\]
Thus a candidate pays only for blockers not already removed. Its
marginal gain can increase as the selected family grows, which is why
screening candidates by their individual insertion costs can discard
profitable exchanges. The final length-33 witness inserts ten supports
with nine distinct blockers, despite twenty blocker incidences.

\subsection{When the candidate representation changes}
A finite candidate pool fixes the coordinates of every candidate.
Its conflict graph therefore describes those coordinates. A
coordinate change can alter an edge, remove a blocker, or admit
a previously excluded point. The geometric lifting criterion
\eqref{eq:lift-mixed} gives the conditions to rebuild the graph
after such a change.

The motion in Section~\ref{sec:motion25}, the signed replacements in
Section~\ref{sec:signed}, and the added shell in dimension 38 change
this finite problem rather than merely improve its incumbent solution.
The two optimizations in dimensions 49--55 also interact: image families
chosen for one line class need not remain optimal after the class grows.
Equation~\eqref{eq:p48-gain} separates the resulting gains.

The moment calculations use a different compression. They retain
all admissible projection signatures and optimize or identify a
linear statistic of their multiplicities. Exact dual multipliers
then express that statistic in terms of known moments. For
dimension 43 the target is determined by an equality; for dimension
45 a nonnegative slack gives a lower bound. The proof requires
completeness of the signature domain, but not prior knowledge of
each individual fibre.

\subsection{Further structural questions}
The constructions leave several precise questions. Which tail Gram
matrices permit the equality case in the reuse bound while supporting
larger head classes? Which boundary conditions allow a signed local
code to attain Proposition~\ref{prop:sign-capacity}? For embedded
sections, which parent root systems force a child count independently
of the ambient lattice, as in Theorem~\ref{d43:thm:structure}?
The last question is distinct from optimizing the determinant of the
section: the two isometric $D_5$ sections in Section~\ref{sec:sections}
already have different perpendicular shell populations.
